% AIrxiv conceptual note (repository-local)
\documentclass[11pt]{article}

\usepackage[margin=1in]{geometry}
\usepackage[T1]{fontenc}
\usepackage[utf8]{inputenc}
\usepackage{lmodern}
\usepackage{microtype}

\usepackage{amsmath,amssymb,amsthm,mathtools}
\usepackage{stmaryrd} % \llbracket \rrbracket
\usepackage{enumitem}
\usepackage{hyperref}
\usepackage{cleveref}

\hypersetup{
  colorlinks=true,
  linkcolor=blue,
  citecolor=blue,
  urlcolor=blue
}

% --- Theorem environments
\newtheorem{definition}{Definition}
\newtheorem{theorem}{Theorem}
\newtheorem{lemma}{Lemma}
\newtheorem{remark}{Remark}

% --- Notation macros
\newcommand{\State}{\mathsf{State}}
\newcommand{\Query}{\mathsf{Query}}
\newcommand{\Evidence}{\mathsf{Evidence}}
\newcommand{\View}{\mathsf{View}}
\newcommand{\side}{\ensuremath{\Sigma}}

\newcommand{\revise}{\mathbin{\oplus}}
\newcommand{\ev}{\mathsf{ev}}
\newcommand{\dsep}{\mathsf{dSep}}
\newcommand{\CI}{\perp\!\!\!\perp}

\title{D-Separation as Certified Rewrite:\\
Connecting Bayesian-Network Semantics to a World-Model Calculus\\
\large (AIrxiv draft/outline; 2026-02-12)}
\author{}
\date{2026-02-12}

\begin{document}
\maketitle

\begin{abstract}
This note records a conceptual and formal ``bridge'' developed while mechanizing d-separation soundness in Lean
and integrating it with a world-model calculus (WM-calculus) for probabilistic reasoning.
The central message is simple: \emph{d-separation is not merely a graph criterion for independence; it is a
\textbf{certificate} that discharges side conditions of compiled probabilistic rewrites}.
In the WM-calculus, inference rules are treated as \side-guarded equalities on query semantics
$e=\ev(W,q)$ rather than as globally valid manipulations of low-dimensional ``truth values.'' The d-separation theorem
becomes the canonical mechanism for turning structure into admissible rewrites.

We highlight three core accomplishments/ideas:
(i) a decomposition of d-separation soundness into small, reusable lemmas (graph layer $\Rightarrow$ semantic layer),
(ii) a discrete proof path that avoids heavy conditional-expectation machinery by reducing key identities to finite sums
and rank-1 factorization on parent fibers, and (iii) a proof-architecture lesson: stateful invariants are essential to
prevent branch explosion when converting between moral-ancestral trails and active trails.

A short philosophical interlude connects the ``truth object + view'' stance to empiricism and phenomenological
interpretations of epistemic state, and notes a resonance with classical theories that treat inference rules as
valid only under explicit preconditions. The note is intended as a primer for future work in the repository.
\paragraph{Status note.} Implementation-status and count claims in this draft are snapshot-scoped; verify against current repository trackers and build outputs before citing as current state. Lean-proven claims are only those tied to explicit Lean file/theorem references; broader mathematical or historical claims are literature-backed context and not asserted as Lean-proven unless explicitly stated.
\end{abstract}

\tableofcontents

\section{Motivation: why d-separation matters for a WM-calculus}

Many neuro-symbolic systems use fast local rules that look like
``truth-value algebra'' (e.g.\ update formulas on $(\text{strength},\text{confidence})$ pairs).
Such rules can be useful, but they are typically \emph{not} globally valid without structural assumptions
about independence and factorization. The WM-calculus stance is:

\begin{quote}
\emph{Make the semantic object primary: for a world-model state $W$ and query $q$, define a truth object
$e=\ev(W,q)$. Treat classical ``truth values'' as views/projections of $e$. Compile fast rules as \side-guarded rewrites,
where \side explicitly states the assumptions under which the rewrite is exact.}
\end{quote}

In this framing, d-separation is an unusually clean bridge:
it turns \emph{graph structure} into \emph{semantic independence facts}, which in turn discharge \side for
many important probabilistic rewrites.

\section{The WM-calculus interface (minimal form)}

\begin{definition}[World-model interface]
A world-model class provides:
\begin{enumerate}[leftmargin=2em]
\item a type of states $\State$ (belief/posterior representation),
\item a type of queries $\Query$,
\item a truth-object type $\Evidence$,
\item an evidence extraction map $\ev : \State \times \Query \to \Evidence$,
\item an update/revision operation $\revise$ on states (possibly partial, to model overlap-aware fusion).
\end{enumerate}
\end{definition}

\begin{definition}[\side-guarded rewrite]
A compiled rule is represented as a theorem schema
\[
  \side \Rightarrow \ev(W,q) \equiv \ev(W,q').
\]
Operationally, the system may apply the rewrite only when \side is discharged (proved) in the current context.
\end{definition}

\begin{remark}[Truth objects vs.\ views]
A view is any map $v : \Evidence \to U$ into a tractable space $U$.
The WM-calculus encourages using multiple views---for speed, control, or interpretability---without confusing them
with the semantics.
\end{remark}

\section{The d-separation bridge as a certified rewrite pipeline}

This section records the conceptual pipeline that guided the mechanization.

\subsection{Graph layer: active trails and moral-ancestral separation}

Classically, d-separation is defined in terms of \emph{active trails} in a DAG.
The standard algorithmic equivalence states that d-separation can be checked by separation in an undirected graph:

\begin{theorem}[Pearl-style equivalence (informal)]
Let $G$ be a DAG, and let $X,Y,Z$ be node sets.
Then $X$ and $Y$ are d-separated by $Z$ in $G$ iff $X$ and $Y$ are separated by $Z$ in the \emph{moralized ancestral
graph} of $X\cup Y\cup Z$.
\end{theorem}

In a mechanized development, this equivalence is best proven by two explicit transforms:
\begin{enumerate}[leftmargin=2em]
\item \textbf{Forward transform:} an active trail $\Rightarrow$ an undirected moral-ancestral trail avoiding $Z$.
\item \textbf{Reverse transform:} a moral-ancestral trail avoiding $Z$ $\Rightarrow$ an active trail.
\end{enumerate}

\paragraph{Mechanization lesson.}
The reverse transform is where naive case splits tend to explode.
A robust proof architecture uses a \emph{stateful invariant} during trail-walking:
``there exists some anchor node in $X$ already actively connected to the current head, and the head is passable.''
This lets the proof proceed step-by-step along the moral trail, updating the anchor when a detour yields a new $X$ node,
and finishing immediately when a detour yields a $Y$ node.

\subsection{Semantic layer: factorization \texorpdfstring{$\Rightarrow$}{->} conditional independence}

The semantic half of d-separation soundness is the global Markov property:
if $\mu$ factorizes according to $G$ (Bayesian network semantics),
then d-separation entails conditional independence:
\[
  \dsep_G(X,Y\mid Z) \Rightarrow X \CI Y \mid Z \quad \text{under } \mu.
\]

For integration into a WM-calculus, the useful payload is not only the independence statement,
but the algebraic identities it implies (``screening-off'') that match compiled rewrite rules.

\subsection{From conditional independence to rewrite equalities}

A typical compiled rewrite needs an identity like
\[
  P(A,C,B)\,P(B) = P(A,B)\,P(C,B),
\]
or equivalently $P(C\mid A,B)=P(C\mid B)$ under positivity assumptions.
In the WM-calculus this becomes a certified rewrite:
\[
  \side_{\text{d-sep}} \Rightarrow \ev(W, q_{\text{lhs}}) = \ev(W, q_{\text{rhs}}),
\]
where $\side_{\text{d-sep}}$ is discharged by a d-separation proof in the graph.

\section{Major formalization breakthroughs (what actually made it work)}

This section collects the key proof-techniques that turned ``standard textbook math'' into Lean-checkable theorems.

\subsection{Breakthrough A: replace heavy \texttt{condExp} reasoning with discrete finite-sum algebra}

A recurring obstacle in measure-theoretic formalizations is conditional expectation.
In the discrete BN setting, the same identities can be proven \emph{more robustly} by finite-sum computation:

\begin{enumerate}[leftmargin=2em]
\item Express $\mu(S)$ as a finite sum of joint weights:
\[
  \mu(S) = \sum_{x} \mathbf{1}_{x\in S}\,w(x)
\]
for measurable $S$.
\item Reduce screening-off on a parent fiber to a rank-1 factorization:
on $\{\mathrm{Pa}(v)=c\}$, sums over descendants collapse by normalization (a telescoping argument),
leaving dependence on $x_v$ only through the local CPT term.
\item Conclude the cross-multiplication identity by algebra on sums.
\end{enumerate}

\paragraph{Takeaway.}
In a discrete mechanization, it is often better to \emph{prove the measure identity directly}
than to force a large conditional-expectation lemma. The former is smaller, faster to elaborate,
and aligns with the actual inference engine used operationally (variable elimination).

\subsection{Breakthrough B: parent-fiber screening via rank-1 factorization}

The key lemma shape is:
for fixed parent assignment $c$ and a set $T$ measurable w.r.t.\ non-descendants excluding parents/self,
\[
  \mu(T \cap \{X_v=a\} \cap \{\mathrm{Pa}(v)=c\}) = \theta(a,c)\cdot K(T,c)
\]
for some $\theta$ determined by the local CPT at $v$ and some $K$ independent of $a$.
Once this is established:
\[
  \mu(T \cap \{X_v\in A\}\cap F)\,\mu(F) = \mu(\{X_v\in A\}\cap F)\,\mu(T\cap F)
\]
follows by summing $\theta(a,c)$ over $a\in A$ and using $\sum_a \theta(a,c)=1$.

\subsection{Breakthrough C: stop branch explosion by carrying state in the reverse transform}

The difficult direction in the moral-ancestral equivalence is:
\[
  \text{moral trail avoiding } Z \Rightarrow \text{active trail}.
\]
A clean approach does \emph{not} demand that each intermediate trail head lies in $X$.
Instead it carries an invariant such as:
\[
  \exists x_0\in X,\ \text{HasActiveTrail}(x_0,\text{head}) \quad \wedge \quad \text{HeadPassable}(\text{head}).
\]
Then each moral edge step either:
\begin{itemize}[leftmargin=2em]
\item extends the state across a direct moral edge,
\item passes through an activated collider witness (if the spouse witness is in $\mathrm{Anc}(Z)$),
\item or triggers a detour lemma: ``if the spouse witness is not in $\mathrm{Anc}(Z)$, then it reaches $X$ or $Y$,''
which either finishes (hit $Y$) or updates the anchor (hit $X$).
\end{itemize}

\paragraph{Why this matters.}
This state-carrying style turns a combinatorial proof into a small number of reusable ``step'' lemmas and one recursion,
which is exactly what proof assistants reward.

\section{How this connects back to WM-calculus and PLN-style inference}

\subsection{D-separation as a \texorpdfstring{$\side$}{Sigma}-discharge engine}

The WM-calculus encourages treating fast rules as \side-guarded rewrites.
In practice, \side often needs to express conditional independence.

Bayesian networks are the canonical class where independence is \emph{structural}:
it can be proven from the graph. Hence:
\[
  \side_{\text{graph}} \Rightarrow \side_{\text{semantic}} \Rightarrow \text{rewrite correctness}.
\]
D-separation is the engine that supplies $\side_{\text{graph}}$.

\subsection{Truth values as views: why local propagation is incomplete}

A recurring formal theme is that correlations live in the joint state.
Therefore no fixed low-dimensional truth value (e.g.\ a pair $(s,c)$) can be globally complete across all queries.
The WM-calculus response is to:
\begin{enumerate}[leftmargin=2em]
\item keep a truth object $e=\ev(W,q)$ that may be structured,
\item expose tractable views for control,
\item use certified \side-guarded rewrites (e.g.\ from d-separation) when a view-level rule is exact.
\end{enumerate}

\subsection{Evidence overlap and partial fusion (why \side matters even outside BNs)}

When evidence overlap is possible, revision should be partial unless disjointness/independence is certified.
This can be modeled by a separation-algebra-like structure for evidence, and by treating ``combine'' steps as admissible
only under explicit preconditions. The same design discipline applies:
\emph{sound speed requires explicit \side.}

\section{Philosophical interlude: empiricism, phenomenology, and rule admissibility}

The WM-calculus is compatible with a broadly empirical reading:
the system does not claim access to ``truth'' as a metaphysical primitive;
it maintains an epistemic state and extracts query-specific evidence objects.
Views (like scalar probabilities or strength/confidence) are tools for action and communication.

Two resonances are worth noting:
\begin{itemize}[leftmargin=2em]
\item \textbf{Phenomenological stance:} a query is an intentional act; the ``truth object'' is the structured content
made available under a particular perspective/intent (here, the current world-model state).
\item \textbf{Rule admissibility tradition:} many classical principles are valid only under explicit preconditions
(e.g.\ decidability, totality, or a chosen semantics). This parallels treating probabilistic rules as admissible rewrites
only under \side.
\end{itemize}

Classical epistemological systems that foreground \emph{conditions of valid inference} align well with the WM-calculus
discipline: inference is licensed by explicit warrants rather than assumed universally.

\section{Outline of the mechanized pipeline (as a roadmap)}

For future work (or new contributors), the mechanization pipeline is:

\begin{enumerate}[leftmargin=2em]
\item \textbf{Graph:} formalize active trails and d-separation; formalize moralization/ancestral construction.
\item \textbf{Equivalence:} prove the moral-ancestral equivalence by forward and reverse transforms using stateful invariants.
\item \textbf{Semantics:} define discrete CPT semantics; expose $\mu$ as a measure and $w(x)$ as joint weights.
\item \textbf{Local Markov:} prove screening-off on parent fibers by finite sums and telescoping elimination.
\item \textbf{Global Markov:} compose graph equivalence with semantic CI to obtain d-separation soundness.
\item \textbf{WM integration:} expose d-separation proofs as \side certificates discharging rewrite rules.
\end{enumerate}

\section{Relationship to companion repository notes}

This note intentionally focuses on the d-separation bridge and the mechanization lessons that make it tractable.
Companion notes cover broader or orthogonal topics:
\begin{itemize}[leftmargin=2em]
\item a general WM-calculus primer (truth objects, views, evidence algebras),
\item algorithmic probability and Solomonoff-style induction notes,
\item process-calculus formalization notes,
\item bootstrapping neurosymbolic agent execution/inference-control notes.
\end{itemize}

\section{Conclusion}

The d-separation theorem is more than a classical result in graphical models:
it is a reusable \emph{certificate generator} for \side-guarded probabilistic rewrites.
Mechanizing it clarifies where independence assumptions enter, how they can be discharged structurally,
and how to keep fast inference rules honest inside a world-model calculus.

\bibliographystyle{plain}
\begin{thebibliography}{9}

\bibitem{Pearl}
J.~Pearl.
\newblock \emph{Probabilistic Reasoning in Intelligent Systems: Networks of Plausible Inference}.
\newblock Morgan Kaufmann, 1988.

\bibitem{KollerFriedman}
D.~Koller and N.~Friedman.
\newblock \emph{Probabilistic Graphical Models: Principles and Techniques}.
\newblock MIT Press, 2009.

\bibitem{Lauritzen}
S.~Lauritzen.
\newblock \emph{Graphical Models}.
\newblock Oxford University Press, 1996.

\end{thebibliography}

\end{document}
